\documentclass[11pt]{article}
\usepackage[utf8]{inputenc}
\usepackage[margin=1in]{geometry}
\usepackage{amsmath,amssymb}
\usepackage{booktabs}
\usepackage{graphicx}
\usepackage{xcolor}
\usepackage[colorlinks=true,linkcolor=blue,citecolor=blue,urlcolor=blue]{hyperref}
\usepackage{authblk}

% ============================================================================
% MAIN CONFERENCE PAPER (full combined study)
% Real measurements: OpenDriveVLA-0.5B; nuScenes-mini open-loop; NeuroNCAP
% closed-loop scene-0103 frontal (10 paired seeds). Honest limitations included.
% ============================================================================

\title{\textbf{Closed-Loop Reveals What Open-Loop Hides: Generation-Coherence\\
Collapse and Group-Wise Recovery in Quantized Autonomous-Driving VLAs}}
\author[1]{Justin Williams}
\author[1]{Kishor Datta Gupta}
\author[1]{Roy George}
\author[1]{Mrinmoy Sarkar}
\affil[1]{Clark Atlanta University \quad \texttt{justinwilliamstech693@gmail.com}}
\date{}

\begin{document}
\maketitle

\begin{abstract}
Post-training quantization of autonomous-driving vision--language--action (VLA) models is widely
reported lossless on open-loop nuScenes L2. We show this is a metric artifact and that the truth
only appears closed-loop. On OpenDriveVLA-0.5B (Qwen2.5-0.5B planner on UniAD perception) we
establish three results. (1) \textbf{Open-loop is blind}: INT8/INT4 weight quantization is lossless
($0.33$\,m), but zeroing the ego-state input inflates L2 to $6.38$\,m ($19\times$)---open-loop L2 is
$\approx$ ego extrapolation and cannot see perception/planning or quantization damage. (2)
\textbf{Closed-loop reveals a failure}: in a real-camera neural-rendering simulator (NeuroNCAP,
two scenes), \emph{naive per-channel} INT4 collapses generation---degenerate trajectory text on
nearly every frame, planner frozen ($0.0\%$ and $8.6\%$ well-formed)---while FP16 stays coherent
($74.3\%$/$72.4\%$). (3) \textbf{Granularity is the fix}: \emph{group-wise} INT4 (group 128,
$\approx 0.1$ extra bits/weight) recovers most of FP16's coherence ($72.9\%$/$57.9\%$). The
collapse is a quantization\,$\times$\,deployment-domain-shift interaction (per-channel INT4 is
lossless in-distribution), governed by scale granularity, not bit-width. We argue
generation-coherence rate is the correct metric (collision counts are a confounded proxy) and that
group-wise INT4 is the actionable, nearly-free recommendation. We are explicit about scope: results
are two-scene with an out-of-distribution renderer (one scene is failed even by FP16), establishing
the mechanism rather than a fleet-wide guarantee.
\end{abstract}

\section{Introduction}
On-board inference motivates quantizing driving VLAs, and 4-bit weight quantization is generally
reported free on open-loop nuScenes planning. We show that conclusion cannot be drawn from
open-loop L2, and that closed-loop evaluation exposes a real, fixable failure. Contributions:
\begin{itemize}
\item A controlled demonstration that open-loop nuScenes L2 is $\approx$ ego-motion extrapolation
($19\times$ ego-ablation), hence structurally blind to quantization damage.
\item The first closed-loop characterization of weight-quantization effects in a driving VLA,
showing per-channel INT4 collapses generation under deployment-domain shift while group-wise INT4
does not.
\item Generation-coherence rate as the un-confounded metric, and a near-free deployment
recommendation (group-wise INT4), with component memory benchmarks toward edge deployment.
\end{itemize}

\section{Background}
\textbf{OpenDriveVLA} pairs a Qwen2.5 language/action head with UniAD stage-1 track+map perception
(a ResNet-101 multi-camera backbone + temporal BEVFormer) as a vision tower; the planner emits a
trajectory as text. \textbf{Open-loop nuScenes} evaluation reports L2 displacement against logged
futures; prior work (``ego status is all you need'') noted it is dominated by ego state.
\textbf{NeuroNCAP} is a photorealistic closed-loop simulator: a NeuRAD neural renderer re-renders
sensor views at the ego's driven pose, so the logged trajectory cannot be coasted on.
\textbf{Group-wise quantization} assigns a scale per block of weights rather than per channel; it
is standard for INT4 LLMs because it limits the damage of weight outliers on the shared scale.

\section{Method}
We apply weight-only post-training quantization (round-to-nearest, dequantized to the compute
dtype) to the 168 Linear layers of the Qwen-0.5B backbone (357.8\,M params); perception,
projectors, embeddings, and head remain FP, isolating the language head's W$b$ effect. We compare
per-output-channel scales against group-wise scales (group sizes 128, 64). For closed-loop we built
a three-process native harness (no Docker): the NeuroNCAP orchestrator, the NeuRAD renderer, and an
OpenDriveVLA model node exposing inference plus an in-place re-quantization endpoint, so techniques
are swept on the loaded model without reloading. The model node reconstructs the UniAD perception
input (lidar2img, CAN-bus, ego-state prompt) from rendered frames each step. We report
\textbf{generation coherence}: the fraction of frames on which the planner emits a well-formed
(non-degenerate) plan.

\section{Results}
\subsection{Open-loop is blind}
\begin{table}[h]\centering
\caption{Open-loop nuScenes-mini ($n{=}162$). L2 in meters (ST-P3 protocol).}
\begin{tabular}{lcc}
\toprule
Setting & L2 (m) & well-formed \\
\midrule
FP16 & 0.33 & 100\% \\
INT8 / INT4 (per-channel) & 0.33 & 100\% \\
INT3 & 1.91 & 92\% \\
INT2 & 6.45 & 0\% \\
FP16, ego-state+history zeroed & \textbf{6.38} ($19\times$) & 100\% \\
\bottomrule
\end{tabular}
\end{table}
INT4 is lossless; the $19\times$ ego-ablation shows why the metric cannot reflect the perception/
planning pathway or its quantization.

\subsection{Closed-loop generation coherence}
\begin{table}[h]\centering
\caption{NeuroNCAP closed-loop, \emph{real-camera} NeuRAD rendering (tcnn), 10 paired seeds per
scene. Generation coherence = well-formed predicted-frame rate. NCAP score / collisions are
secondary (confounded) proxies.}
\begin{tabular}{lcccccc}
\toprule
 & \multicolumn{3}{c}{0103 (frontal)} & \multicolumn{3}{c}{0796 (stationary)} \\
\cmidrule(lr){2-4}\cmidrule(lr){5-7}
Config & coher. & NCAP & coll. & coher. & NCAP & coll. \\
\midrule
FP16              & 74.3\%          & 5.0  & 0/10  & 72.4\%          & 0.88 & 10/10 \\
INT4 per-channel  & \textbf{0.0\%}  & 4.08 & 4/10  & \textbf{8.6\%}  & 0.00 & 10/10 \\
INT4 group-128    & \textbf{72.9\%} & 5.0  & 0/10  & \textbf{57.9\%} & 1.58 & 10/10 \\
\bottomrule
\end{tabular}
\end{table}
On both scenes per-channel INT4 collapses generation ($0.0\%$/$8.6\%$ well-formed) while group-wise
INT4 recovers most of FP16's coherence ($72.9\%$/$57.9\%$ vs $74.3\%$/$72.4\%$). The \emph{same}
per-channel INT4 is lossless in-distribution (open-loop, Table 1), so the failure is a
quantization\,$\times$\,domain-shift interaction. Collision/NCAP is confounded and scene-dependent:
on 0796 \emph{every} config including FP16 collides $10/10$ (the base policy fails that scenario),
so only coherence cleanly separates the configs there; on 0103 collisions do track collapse
($4/10$ vs $0/10$). We therefore lead with coherence, and we do not claim group-wise exactly equals
FP16 (it matches on 0103, partially recovers on 0796).

\subsection{Why granularity}
A shared INT4 scale $s=\max|w|/7$ is set by the group's largest weight; one outlier coarsens the
grid for all weights it shares the scale with. Group-wise scales confine each outlier to its block,
preserving precision on the signal-carrying weights. Under distribution shift the planner's hidden
states drift; the per-channel coarsened grid then tips greedy decoding into malformed trajectory
text, whereas the group-wise grid preserves coherence---at a cost of $16/128=0.125$ bits/weight.

\subsection{Memory benchmarks (toward edge)}
\begin{table}[h]\centering
\caption{Backbone weight footprint and overhead (357.8\,M Linear params); effective bits include
per-group FP16 scales.}
\begin{tabular}{lcc}
\toprule
Precision & eff. bits/weight & backbone weights \\
\midrule
FP16 & 16.0 & $\approx 1.0$\,GB \\
INT4 per-channel & 4.06 & $\approx 0.25$\,GB \\
INT4 group-128 & 4.125 & $\approx 0.26$\,GB \\
\bottomrule
\end{tabular}
\end{table}
The robust recipe (group-wise INT4) is $\approx 0.26$\,GB; the full inference node fits in
$3$--$5$\,GB. Perception (FP32, custom CUDA ops) dominates wall-clock, so the quantization win is
robustness, not latency; the edge bottleneck is perception, not the planner (see companion systems
study).

\section{Discussion}
The two halves compose into one message: open-loop L2 \emph{hides} quantization damage because it is
ego-extrapolation, and closed-loop \emph{reveals} it because perception/planning are on the critical
path. The damage, when revealed, is a granularity artifact: per-channel INT4 is fine in-distribution
but loses generation coherence under deployment shift, while group-wise INT4 is robust. Practitioners
should (i) never certify a quantized driving policy on open-loop L2, and (ii) default to group-wise
INT4 for the planner.

\section{Limitations and threats to validity}
(1) \textbf{Two scenes} (NeuroNCAP 0103 frontal, 0796 stationary); the mechanism is established and
replicated once, not fleet-wide. The remaining NeuRAD-reconstructed scenes need gated nuScenes
trainval metadata. (2) \textbf{Out-of-distribution renderer}: rendered inputs are OOD to real
nuScenes---even FP16 fails 0796's collision test ($10/10$)---so absolute NCAP/collision numbers are
not comparable to in-distribution operation and our claims are \emph{relative} (group $\approx$ FP16
$\gg$ per-channel on coherence). (1b) \textbf{Partial recovery}: group-wise matches FP16 coherence
on 0103 but only partially on 0796 ($57.9\%$ vs $72.4\%$). (3) The collision
metric is a confounded downstream proxy (collapse$\to$frozen$\to$geometry-dependent crash, mediated
by the degenerate fallback); we therefore lead with coherence. (4) An activation-aware (AWQ) low-bit
variant is excluded pending a fix. (5) Fake-quant (dequantized) weights are used, so packed memory/
latency gains are projected, not measured on-device.

\section{Conclusion}
For quantized autonomous-driving VLAs, open-loop L2 is blind ($19\times$ ego-ablation) and
closed-loop is necessary: there, naive per-channel INT4 collapses generation under deployment-domain
shift while group-wise INT4 ($\approx 0.1$ extra bits/weight) matches FP16. Use closed-loop
coherence as the metric and group-wise INT4 as the recipe.

\begin{thebibliography}{9}
\bibitem{opendrivevla} OpenDriveVLA. 2025. (Replace with verified citation.)
\bibitem{uniad} Hu et al. \emph{Planning-oriented Autonomous Driving (UniAD)}. CVPR 2023.
\bibitem{egostatus} Zhai et al. \emph{Rethinking Open-Loop Evaluation of End-to-End AD}. 2023.
\bibitem{neuroncap} NeuroNCAP. ECCV 2024.
\bibitem{neurad} NeuRAD: \emph{Neural Rendering for Autonomous Driving}. CVPR 2024.
\bibitem{awq} Lin et al. \emph{AWQ}. MLSys 2024.
\bibitem{gptq} Frantar et al. \emph{GPTQ}. ICLR 2023.
\end{thebibliography}
\end{document}
